\section{Do blind judges know good copy?}
\label{sec:taste}

hill-climb decides taste rounds with a panel of three fresh-context judges who see the two versions under
neutral labels, each with its own coin flip for which side is new, and keeps a change only on a strict
majority. The panel guards against position bias and a single judge's whims. Whether it tracks what an audience
actually prefers is a separate question, and the one that matters.

\paragraph{Data.} The Upworthy Research Archive~\cite{upworthy} holds thousands of real headline A/B tests with
impressions and clicks. From the exploratory release we took every test with two headlines that ran with the
same image, excerpt and lede, kept the pair with the largest click-rate gap in each test when that gap had a
two-proportion $z \ge 3$ and each arm had at least 1{,}000 impressions, and sampled 100 such pairs with a fixed
seed. These are decisive tests: the median winner drew 177\% more clicks than the loser.

\paragraph{Protocol.} Each pair went to three independent judges per panel, using the skill's judge prompt
template filled for headlines (three axes, then a verdict of A, B or TIE), with sides assigned by a recorded
coin flip per judge. We ran three panels, one per model size (Claude Haiku, Sonnet and Opus), 900 calls in all.
No call failed to parse.

\begin{table}[t]\centering\small
\caption{Three-judge panels on 100 decisive Upworthy headline tests. Brackets: 95\% Wilson intervals (bootstrap for $\kappa$). Chance is 50\%.}\label{tab:taste}
\begin{tabular}{lcccc}\toprule
Panel & One judge right & Majority right & Right when unanimous & Fleiss $\kappa$ \\\midrule
Haiku & 63\% [57, 68] & 65\% [55, 74] & 70\% [58, 80] (61) & 0.45 [0.31, 0.58] \\
Sonnet & 69\% [63, 74] & 69\% [59, 77] & 74\% [63, 82] (80) & 0.71 [0.59, 0.82] \\
Opus & 69\% [63, 74] & 67\% [57, 75] & 73\% [63, 81] (88) & 0.83 [0.74, 0.92] \\
\bottomrule\end{tabular}\end{table}


\paragraph{Findings.} Judges agree with each other far more than with readers. Agreement rises with model size,
from a Fleiss $\kappa$ of 0.45 for Haiku to 0.83 for Opus, and the Sonnet and Opus panels return the same majority
verdict on 84\% of pairs. Accuracy against the real winner does not follow: every panel lands between 65\% and
69\%, close to the 65\% ceiling that LOLA~\cite{lola} reported for prompted models on the same archive. A majority of
three is no more accurate than a single judge, because the three judges make the same mistakes. Unanimous
panels are slightly more accurate (70 to 74\%), so unanimity carries a little signal. Position bias is
modest and cancelled by the per-judge flips (Haiku picks slot A 36\% of the time, Opus 52\%).

\paragraph{What this means for the skill.} The panel does its narrow job: it gives a stable, unbiased verdict
on what a capable reader would prefer. It does not measure an audience. A climb that a panel scores can drift
toward what models like. We added this result to the skill's judging guide: treat a panel win as ``judges
preferred'', not ``users will prefer'', use a real outcome measure as the score when one exists, and confirm the
final best against it when it is costly to measure every round.
